\subsection{The remaining unstable algebraic moduli stack}

\begin{result}{E.T6}{Proved}{The unmarked genus-one algebraic stack}
Let $\mathcal M_{1,0}$ denote the stack of smooth proper
genus-one curves over $k=\overline\Q$, without a chosen
origin.  Its group of $k$-linear algebraic autoequivalences,
modulo $2$-isomorphism, is $C_2$.
Writing $\mathcal B=\mathcal M_{1,1}$ and
$\mathcal E\to\mathcal B$ for the universal elliptic curve,
the comparison and the nontrivial automorphism are
\[
 \mathcal M_{1,0}\simeq B_{\mathcal B}\mathcal E,
 \qquad (E,T)\longmapsto(E,T\mathbin{\wedge}^{E}Q_E).
 \tag{E.T9}
\]
Here $Q$ is the nonzero class in
$H^1(\mathcal B,\mathcal E)=C_2$, obtained from the
$\mathcal E[2]$-torsor specified by the cocycle below.
The induced homomorphism
$\Aut_k(\mathcal M_{1,0})\to\Aut_k(\mathcal M_{1,1})$
is trivial.  In particular the quadratic twist (E.T8)
does not lift through the Jacobian morphism, whereas
(E.T9) acts trivially on its base.
\end{result}

\begin{proof}
A genus-one curve is naturally a torsor under its degree-zero
Picard group; after choosing a local point this is the usual
simply transitive elliptic-curve action.  The action is
independent of the point and descends.  Conversely an
elliptic-curve torsor is a smooth proper genus-one curve.
These constructions and their maps are inverse, proving
the first equivalence in (E.T9).  Its map to $\mathcal B$
is also the quotient of inertia by its connected identity
component: the latter is precisely the elliptic group of
translations.  An autoequivalence preserves this connected
component.  In local torsor trivializations it therefore
descends to the Jacobian object and its isomorphisms;
the descended functors agree on overlaps.  Thus every
autoequivalence induces an automorphism of $\mathcal B$.

The nontrivial automorphism (E.T8) of $\mathcal B$ cannot
occur in that image.  At an elliptic curve with $j=1728$
it inverts the order-four stabilizer.  Conjugation by
that stabilizer on the one-dimensional tangent of the
translation group has a faithful character of order four.
An isomorphism between the translation groups would have
an invertible scalar derivative intertwining this character
and its inverse.  The two values on a generator are $i$
and $-i$, so the intertwining equation has no nonzero
solution.  Compatibility of the inertia homomorphism
with conjugation gives exactly this impossible equation.
E.T4 now proves that every automorphism induces the
identity class on $\mathcal B$.

We compute the relative automorphisms and the possible
identifications of the chosen comparison with that identity.
An equivalence of $B_{\mathcal B}\mathcal E$ over
$\mathcal B$ sends its zero object to an $\mathcal E$-torsor
$Q$ and induces an automorphism $\epsilon$ of its band
$\mathcal E$.  Locally trivializing $Q$ identifies the
functor with pushforward of torsors by $\epsilon$;
on overlaps the changes are precisely the translations
defining $Q$.  It is therefore
$T\mapsto Q\mathbin{\wedge}^{\mathcal E}\epsilon_*T$,
and this description also detects all natural
isomorphisms over the base.  The relative group is
\[
 H^1(\mathcal B,\mathcal E)\rtimes\{1,-1\}.
 \tag{E.T10}
\]
The only band automorphisms are $\pm1$: at the generic
Weierstrass curve, both $a$ and $b$ are nonzero and a
group automorphism has scale $u$ with $u^4=u^6=1$,
so $u^2=1$.  Equality at the generic point implies equality
everywhere by separatedness.  The global $2$-automorphisms
of $\operatorname{id}_{\mathcal B}$ are likewise precisely
$\pm1$, by their generic values.  The negative one
identifies $(Q,\epsilon)$ with $(Q,-\epsilon)$ after
the base comparison is forgotten: on a torsor with its
action reversed the identity on the underlying curve
is equivariant for the base elliptic isomorphism $[-1]$.
It sends the torsor class to $-Q$, which equals $Q$
by the calculation below.  Thus the absolute group is
the quotient of (E.T10) by the subgroup generated by
$(0,-1)$.

First, $\mathcal E(\mathcal B)=0$.  A section must be fixed
by the universal involution $[-1]$, hence is a section of
$\mathcal E[2]$.  Its monodromy is the standard
$\operatorname{SL}_2(\F_2)$ action, which fixes only zero.
Every class in $H^1(\mathcal B,\mathcal E)$ is torsion.
To prove this, pull it back to a connected finite etale
level cover which is a smooth affine curve.  A torsor
under its elliptic scheme is proper.  A generic section
extends over every discrete valuation ring by properness,
and uniqueness makes these extensions a global section.
Consequently restriction to the function field is
injective.  Over that field a torsor obtains a point
over a finite separable extension; restriction followed
by the norm shows that its class is killed by the
degree.  It is therefore torsion on the level curve.
For the finite etale level cover, sum over the sheets
defines the norm on the abelian sheaf and its torsors;
norm after restriction is multiplication by its degree.
This proves torsion on $\mathcal B$ as well.

The Kummer sequence, the zero section calculation, and
finite-cover comparison now identify
\[
 H^1(\mathcal B,\mathcal E)[n]
   =H^1(\mathcal B,\mathcal E[n])
   =H^1\bigl(\operatorname{SL}_2(\Z),(\Z/n)^2\bigr).
 \tag{E.T11}
\]
For the last equality, finite torsors with the standard
linear action are exactly affine monodromy actions of
the orbifold group; their choices of a sheet differ
by coboundaries.  This proves the required degree-one
comparison directly from Riemann existence.

Put $M=(\Z/n)^2$ and use
\[
 R=\begin{pmatrix}0&-1\\1&0\end{pmatrix},\qquad
 S=\begin{pmatrix}0&-1\\1&1\end{pmatrix},\qquad
 \operatorname{SL}_2(\Z)=C_4*_{C_2}C_6,
 \quad R^2=S^3=-I.
\]
The amalgam description follows also from the braid
presentation in B.9.  The group-cohomology exact sequence
of its tree is obtained by applying equivariant cochains
to its two vertex orbits and its one edge orbit.  In
degrees zero and one it contains
\[
 \begin{split}
 M^{C_4}\oplus M^{C_6}\longrightarrow M^{C_2}
 &\longrightarrow H^1(\operatorname{SL}_2(\Z),M)\\
 &\longrightarrow H^1(C_4,M)\oplus H^1(C_6,M)
 \longrightarrow H^1(C_2,M).
 \end{split}
 \tag{E.T12}
\]
Here $S-I$ has determinant one, so both terms for $C_6$
vanish.  Also $M^{C_2}=M[2]$ and
$M^{C_4}=\ker(R-I)$.  For odd $n$ both are zero.
For even $n$ they have orders four and two, respectively.
The norm $1+R+R^2+R^3$ is zero; hence
$H^1(C_4,M)=M/(R-I)M$ is zero for odd $n$ and $C_2$
for even $n$.  The restriction to
$H^1(C_2,M)=M/2M$ sends the generator represented by
$(1,0)$ to $(1+R)(1,0)=(1,1)$ modulo $2M$.
It is injective.  Exactness proves
\[
 H^1(\operatorname{SL}_2(\Z),M)=
 \begin{cases}0& n\text{ odd},\\ C_2& n\text{ even}.
 \end{cases}
 \tag{E.T13}
\]
The inclusion $\mathcal E[n]\hookrightarrow\mathcal E[2n]$
is multiplication by two in these coordinates.  On
$M[2]/\ker(R-I)$ it carries the nonzero class represented
by $(n/2,0)$ to that represented by $(n,0)$.  Thus the
even groups in (E.T13) have compatible nonzero classes.
More directly, torsion and (E.T11)--(E.T13) imply that
every class has order at most two, and its two-torsion
group is $C_2$.  This proves
$H^1(\mathcal B,\mathcal E)=C_2$ without a restriction
on the odd parts of the torsion orders.

An explicit cocycle for $n=2$ is
\[
 c=(1,0),\qquad z(R)=0,\qquad z(S)=(S-I)c=(1,1).
 \tag{E.T14}
\]
The norm relations hold and the two values on
$R^2=S^3$ are zero.  A coboundary with $z(R)=0$
would come from an $R$-fixed vector, whereas
$z(S)=(S-I)c$ forces that vector to equal $c$;
$Rc\ne c$.  This proves its nontriviality.
Its affine permutations of $\F_2^2$ generate $S_4$
and act transitively on four sheets.  The exact program
\path{code/e_unmarked_genus_one.py} checks this assertion,
the amalgam cocycle and coboundary counts for
$1\leq n\leq16$ and $n=24,32,64$, and the maps to $2n$.
It was run; its JSON certificate and complete output
are included.  Algebraizing this finite torsor and
extending its structure group gives $Q$ in (E.T9).
Its class has order two, so twisting by it is an
involutive autoequivalence.  Evaluating on the zero
torsor detects its nontriviality even after forgetting
the base comparison.  Quotienting (E.T10) as above
proves the stated full group.
\end{proof}
